\section{Background}
\subsection{Transformer Self-Attention}
Let $X \in \mathbb{R}^{N \times d_{\text{model}}}$ be an input sequence of $N$ tokens, each embedded into a $d_{\text{model}}$-dimensional vector. The self-attention mechanism transforms $X$ into queries $Q$, keys $K$, and values $V$ via learned linear projections:
\[
Q = X W_Q,\quad K = X W_K,\quad V = X W_V,
\]
where $W_Q, W_K, W_V \in \mathbb{R}^{d_{\text{model}} \times d_k}$ and typically $d_k = d_{\text{model}} / h$ for $h$ heads. The attention output is:
\[
\text{Attention}(Q,K,V) = \text{softmax}\left( \frac{Q K^\top}{\sqrt{d_k}} \right) V.
\]
The $QK^\top$ matrix is $N \times N$, leading to $O(N^2 d_k)$ time and $O(N^2)$ memory per head. In multi-head attention, these computations are performed in parallel across $h$ heads and then concatenated.

\subsubsection{Quadratic Complexity: A Deeper Look}
The quadratic bottleneck arises because each query must attend to all keys. For $N=512$, $N^2 = 262,144$; for $N=4096$, $N^2 = 16.7$ million—a 64-fold increase. In practice, this limits BERT to sequences of at most 512 tokens. For tasks like document classification, question answering over long passages, or speech recognition, this is insufficient.

\subsection{Singular Value Decomposition (SVD)}
For any matrix $H \in \mathbb{R}^{N \times d}$ (assuming $N \ge d$ without loss of generality), the full SVD is:
\[
H = U \Sigma V^\top,
\]
where $U \in \mathbb{R}^{N \times N}$ is orthogonal, $V \in \mathbb{R}^{d \times d}$ is orthogonal, and $\Sigma \in \mathbb{R}^{N \times d}$ is a diagonal matrix with singular values $\sigma_1 \ge \sigma_2 \ge \dots \ge \sigma_d \ge 0$ on the diagonal. The columns of $U$ are the left singular vectors, and the columns of $V$ are the right singular vectors.

\subsubsection{Truncated SVD}
For a given rank $r < \min(N,d)$, the truncated SVD keeps only the first $r$ singular values and vectors:
\[
H_r = U_r \Sigma_r V_r^\top,
\]
where $U_r \in \mathbb{R}^{N \times r}$, $\Sigma_r \in \mathbb{R}^{r \times r}$, $V_r \in \mathbb{R}^{d \times r}$. The Eckart–Young theorem states that $H_r$ is the best rank-$r$ approximation to $H$ in both the Frobenius norm and the spectral norm:
\[
\| H - H_r \|_F = \sqrt{\sum_{i=r+1}^{\min(N,d)} \sigma_i^2}, \quad
\| H - H_r \|_2 = \sigma_{r+1}.
\]

\subsubsection{Randomized SVD}
Computing the full SVD of a large matrix costs $O(N d \min(N,d))$, which is too expensive for every forward pass. However, we only need the top $r$ left singular vectors. The randomized SVD algorithm \cite{halko2011finding} computes an approximate rank-$r$ SVD in $O(N d \log r + (N+d) r^2)$ time, which is linear in $N$ and $d$. The algorithm proceeds as:
\begin{enumerate}
    \item Form $Y = H \Omega$, where $\Omega \in \mathbb{R}^{d \times (r+p)}$ is a random Gaussian matrix ($p$ is oversampling, typically 10).
    \item Compute orthonormal basis $Q$ via QR decomposition of $Y$.
    \item Project $H$ onto $Q$: $B = Q^\top H$.
    \item Compute SVD of $B$: $B = \hat{U} \Sigma V^\top$.
    \item Then $U = Q \hat{U}$ gives the approximate left singular vectors.
\end{enumerate}
This algorithm is highly efficient and accurate when the singular values decay moderately.

\subsection{Low-Rank Attention: Prior Work}
\subsubsection{Linformer}
Linformer \cite{wang2020linformer} proposes to replace the $N \times N$ attention matrix with $N \times r$ by projecting keys and values using fixed matrices $E$ and $F$:
\[
\tilde{K} = E^\top K,\quad \tilde{V} = F^\top V,
\]
where $E, F \in \mathbb{R}^{N \times r}$ are learned or random. The attention becomes:
\[
\text{Attention}(Q,K,V) \approx \text{softmax}\left( \frac{Q (E^\top K)^\top}{\sqrt{d_k}} \right) (F^\top V).
\]
Complexity reduces to $O(N r d)$. However, $E$ and $F$ are independent of the input, which may not be optimal.

\subsubsection{Performer}
Performer \cite{choromanski2020rethinking} uses kernel methods to approximate softmax with random features, achieving $O(N d^2)$ complexity without an explicit low-rank assumption. While theoretically appealing, the approximation error can be high for certain tasks.

\subsubsection{Nyströmformer}
Nyströmformer \cite{xiong2021nystromformer} uses the Nyström method to approximate the attention matrix by sampling landmark points. It also achieves linear complexity but requires selecting landmarks, which may be suboptimal.

Our SVD-Lin differs by using the input itself to derive the projection, making it adaptive and theoretically optimal for the given rank.